% Further consequences; no priority claim. Owned by /root/frobenius.
\paragraph{C.4r (Proved: a single surface-cycle class detects the remaining complex cases).}
Let $A/\mathbb C$ be an abelian variety of dimension $g\in\{2,3\}$,
put $X=A\times\widehat A$, and write
$P=c_1(\mathcal P_A)\in H^2(X,\mathbb Z)$.
The following are equivalent:
\begin{enumerate}
\item $A$ is an actual direct summand of a product of Jacobians;
\item the integral Hodge class
\[
                 \frac{P^{2g-2}}{(2g-2)!}
                 \in H^{4g-4}(X,\mathbb Z)
\]
is the class of an integral algebraic cycle.
\end{enumerate}
Thus the remaining surface case is exactly the algebraicity of
$P^2/2$ on the fourfold $A\times\widehat A$, and the remaining
threefold case is exactly that of $P^4/24$ on the sixfold
$A\times\widehat A$. In both cases the requested cycles have
dimension two.

\emph{Proof.}
If (1) holds, \citet[Theorem 1.4]{cyVoisin2024} and
\citet[Theorem 1.1]{cyBeckmannFortman2023} imply that every
component of $\operatorname{ch}(\mathcal P_A)$ is integrally
algebraic, proving (2). These divided powers are integral in
cohomology; integrality does not mean their indicated rational
intersection expressions are already integral in the Chow group.

Conversely, intersect an integral cycle representing the class
in (2) with the integral divisor class $P$. The resulting integral
curve cycle has class
\[
 \frac{P^{2g-1}}{(2g-2)!}
       =(2g-1)\frac{P^{2g-1}}{(2g-1)!}.
\]
Apply the integral Matsusaka factorization
\eqref{eq:c4-matsusaka-integer} on $X$ with the unimodular
Poincar\'e line bundle. Since its Euler characteristic is
$(-1)^g$, we obtain $(2g-1)\in I_J(X)$, after adjusting the
sign of one homomorphism, and therefore $(2g-1)\in I_J(A)$.
For $g=2$, C.4n already gives $2\in I_J(A)$, so
$1=3-2\in I_J(A)$. For $g=3$, C.4p gives
$6\in I_J(A)$, so $1=6-5\in I_J(A)$.
The explicit splitting of C.4h proves (1).
\hfill$\square$

This is an exact reformulation, not a proof of either remaining
algebraicity assertion. In particular, rational algebraicity of
the cohomological Fourier isometry of an abelian surface does
not supply the integral cycle in (2).

\paragraph{C.4s (Known consequence: every supersingular surface is affirmative).}
Every supersingular abelian surface over an algebraically closed
field of positive characteristic is an actual direct summand of
a product of Jacobians, including in characteristic two.
Indeed, \citet[Lemma 6.2, p.~185]{frC4FuLi2021} prove that it
admits a principal polarization. Their Proposition 6.1,
p.~184, explicitly includes characteristic two in the
N\'eron--Severi lattice calculation used by the proof.
The principally polarized surface case C.4f then applies.
This does not settle arbitrary nonsupersingular surfaces whose
polarization Euler characteristics are all even.
